%!TEX root = README.tex
\documentclass[10pt,a4paper]{article}

\usepackage[margin=0.8in]{geometry}
\usepackage{amsmath,amssymb}
\usepackage{xcolor}
\usepackage[colorlinks=true,linkcolor=blue!75!black,urlcolor=blue!75!black]{hyperref}
\usepackage{parskip}
\usepackage{fancyvrb}
\usepackage{booktabs}
\usepackage{enumitem}

\title{\textbf{University Homework \& Quiz \LaTeX{} Template}\\[0.25em]\large A General-Purpose Bilingual Academic Template for Assignments and Quizzes}
\author{University of Tehran}
\date{}

\begin{document}
\maketitle

\section{Introduction}
This repository provides a versatile, general-purpose \textbf{academic document template} designed for university courses across computer science, engineering, and mathematics. Originally adapted for courses at the \textbf{University of Tehran}, it is fully course-agnostic and provides two specialized document classes:
\begin{itemize}[leftmargin=1.5em, itemsep=0.2em]
    \item \textbf{\texttt{homework.cls}} --- For comprehensive homework assignments, problem sets, and lab handouts with multiple official header styles.
    \item \textbf{\texttt{quiz.cls}} --- For short exams, quizzes, and classroom assessments with student information fields and framed answer boxes.
\end{itemize}

\section{Key Features}
\begin{itemize}[leftmargin=1.5em, itemsep=0.2em]
    \item \textbf{Dual Document Classes:} Specialized classes for both homework assignments and quizzes.
    \item \textbf{Multiple Header Styles (\texttt{homework.cls}):} Switch between Classic Banner (\texttt{header-type=1}) and Modern Wide Side-by-Side (\texttt{header-type=2}) with large 2.0cm logos.
    \item \textbf{Quiz Environment (\texttt{quiz.cls}):} Dedicated header with exam duration (\texttt{\textbackslash quiztime}), student name, student ID fields, and answer boxes (\texttt{\textbackslash answerbox}).
    \item \textbf{Smart Asset \& Font Locator:} Automatically searches up to 5 nested folder levels for \texttt{Fonts/} and \texttt{Background/} assets.
    \item \textbf{Bilingual \& Persian Typesetting (XePersian):} Bundled standard Persian fonts (\textit{XB Niloofar}, \textit{XB Zar}, \textit{IRZar}, \textit{IRLotus}, \textit{XB Yas}).
    \item \textbf{Modular Subfiles Structure:} Organize questions into standalone subfiles combined seamlessly via \texttt{\textbackslash subfile\{...\}}.
    \item \textbf{Conditional Solutions:} Toggle student solutions globally using the \texttt{compileanswers} boolean flag.
    \item \textbf{Syntax-Highlighted Code:} Formats source code elegantly using \texttt{minted} for Python, C/C++, Verilog, Java, etc.
\end{itemize}

\section{Repository Structure}
\begin{itemize}[leftmargin=1.5em, itemsep=0.2em]
    \item \texttt{homework.cls} --- Homework document class supporting multiple header designs and margin extensions.
    \item \texttt{quiz.cls} --- Quiz document class supporting student information fields and answer boxes.
    \item \texttt{Background/} --- High-resolution University of Tehran seal and Faculty of Engineering emblem.
    \item \texttt{Fonts/} --- High-quality Persian TrueType fonts bundled directly in the repository.
    \item \texttt{HW.tex} --- Master document demonstrating homework layout with \texttt{header-type=2}.
    \item \texttt{latexmkrc} --- Automated build configuration for Overleaf and local \texttt{latexmk} XeLaTeX builds.
    \item \texttt{LaTeX-build/} --- Standard build directory for compiler artifacts (\texttt{.pdf}, \texttt{.log}, \texttt{.aux}).
\end{itemize}

\section{Document Class Options}
\begin{itemize}[leftmargin=1.5em, itemsep=0.2em]
    \item \texttt{largemargins} --- Widens page margins to leave ample space for grading annotations (both classes).
    \item \texttt{newpage} --- Automatically begins each new question on a fresh page (both classes).
    \item \texttt{anonymous} --- Omits personal contributor details for blind grading (both classes).
    \item \texttt{header-type=1|2} --- Selects Classic Banner (\texttt{1}) or Modern Wide (\texttt{2}) header in \texttt{homework.cls}.
    \item \texttt{nostudentinfo} --- Suppresses the student name, ID, and duration fields in \texttt{quiz.cls}.
\end{itemize}

\clearpage

\section{Homework Configuration Example (\texttt{homework.cls})}
\begin{Verbatim}[frame=lines,framesep=2mm,fontsize=\footnotesize]
\documentclass[12pt,largemargins,header-type=2]{homework}

\newcommand{\eventclass}{Automata & Language Theory} % Course name
\newcommand{\eventsemester}{Spring 2026}             % Semester
\newcommand{\eventtype}{Homework}                   % Assignment type
\newcommand{\eventnum}{2}                           % Assignment number
\newcommand{\eventinstructor}{Dr. Fadaei & Dr. Yaghoobzadeh} % Instructors
\newcommand{\eventcntrbs}{Mohammad Khoshroo}        % Author or TA
\newcommand{\hwemail}{ta-email@ut.ac.ir}            % Contact email
\newcommand{\eventdatetype}{Due Date:}
\newcommand{\eventDeadlineDay}{8}
\newcommand{\eventDeadlineMonth}{May}
\newcommand{\eventDeadlineYear}{2026}
\newcommand{\eventDeadlineHour}{23:59}

\begin{document}
\makeheader
\subfile{Q-A/Q1.tex}
\clearpage
\subfile{Q-A/Q2.tex}
\end{document}
\end{Verbatim}

\section{Quiz Configuration Example (\texttt{quiz.cls})}
\begin{Verbatim}[frame=lines,framesep=2mm,fontsize=\footnotesize]
\documentclass[12pt,largemargins]{quiz}

\newcommand{\eventclass}{Discrete Mathematics}      % Course name
\newcommand{\eventtype}{Quiz}                       % Exam type
\newcommand{\eventnum}{9}                           % Quiz number
\newcommand{\eventname}{Number Theory}              % Topic
\newcommand{\eventcntrbs}{Instructor: F. Mohammadi} % Instructor
\newcommand{\eventdatetype}{Date:}                  % Date prefix
\newcommand{\eventdate}{2026/05/10}                 % Date
\newcommand{\quiztime}{15 Minutes}                  % Exam duration

\begin{document}
\maketitle

\question
Find all positive integers $a$ and $b$ such that:
\[ \frac{1}{a} + \frac{1}{b} = \frac{1}{101} \]

\answerbox{2.5cm} % Framed answer box for students

\solution
Detailed solution is written here.
\end{document}
\end{Verbatim}

\section{Metadata Variables Reference}
\begin{itemize}[leftmargin=1.5em, itemsep=0.2em]
    \item \texttt{\textbackslash eventclass\{...\}} --- Course title (e.g. Automata Theory, AI).
    \item \texttt{\textbackslash eventtype\{...\}} --- Assignment or quiz type (e.g. Homework, Quiz, Project).
    \item \texttt{\textbackslash eventnum\{...\}} --- Assignment or quiz number.
    \item \texttt{\textbackslash eventname\{...\}} --- Topic or chapter subtitle.
    \item \texttt{\textbackslash eventsemester\{...\}} --- Academic semester (e.g. Spring 2026).
    \item \texttt{\textbackslash eventinstructor\{...\}} --- Instructor name(s).
    \item \texttt{\textbackslash eventinstructortype\{...\}} --- Custom instructor label prefix (default: Instructor).
    \item \texttt{\textbackslash eventcntrbs\{...\}} --- Teaching assistants, designers, or authors.
    \item \texttt{\textbackslash eventcntrbstype\{...\}} --- Custom contributor label prefix (default: Designers).
    \item \texttt{\textbackslash hwemail\{...\}} --- Contact email address (automatically wrapped in a \texttt{mailto:} link).
    \item \texttt{\textbackslash eventdatetype\{...\}} --- Date label prefix (default: Due Date: or Date:).
    \item \texttt{\textbackslash eventdate\{...\}} --- Unified date string.
    \item \texttt{\textbackslash eventDeadlineDay}, \texttt{\textbackslash eventDeadlineMonth}, \texttt{\textbackslash eventDeadlineYear}, \texttt{\textbackslash eventDeadlineHour} --- Granular deadline components.
    \item \texttt{\textbackslash eventcourselink\{...\}} --- URL pointing to the course web page.
    \item \texttt{\textbackslash quiztime\{...\}} --- Quiz duration / allowed response time (e.g. 15 Minutes).
    \item \texttt{\textbackslash quiztimelabel}, \texttt{\textbackslash studentnamelabel}, \texttt{\textbackslash studentidlabel} --- Customizable field labels for quizzes.
\end{itemize}

\section{Question, Solution \& Utility Commands}
\begin{itemize}[leftmargin=1.5em, itemsep=0.2em]
    \item \texttt{\textbackslash makeheader} or \texttt{\textbackslash maketitle} --- Generates the official course header box (both classes).
    \item \texttt{\textbackslash setheadertype\{1|2\}} --- Dynamically switches header style at runtime (\texttt{homework.cls}).
    \item \texttt{\textbackslash question} --- Starts a numbered question with a colored dividing rule.
    \item \texttt{\textbackslash question*\{Title\}} --- Starts a question header with a custom title string.
    \item \texttt{\textbackslash question\#\{Points\}} --- Creates a question header with point allocation badge (\texttt{homework.cls}).
    \item \texttt{\textbackslash solution} --- Creates a solution header with a red dividing rule.
    \item \texttt{\textbackslash solution*\{Title\}} --- Custom solution header.
    \item \texttt{\textbackslash answerbox\{height\}} --- Generates an empty framed box of specified height for student answers (both classes).
    \item \texttt{\textbackslash ExternalLink} --- Renders an inline external link TikZ icon.
    \item \texttt{\textbackslash RN\{number\}} --- Converts an integer to uppercase Roman numerals (e.g. \texttt{\textbackslash RN\{4\}} gives IV).
    \item \texttt{alphaparts} / \texttt{arabicparts} --- Sub-question environments for (a), (b) or 1.1, 1.2 numbering.
    \item \texttt{\textbackslash questiontype} / \texttt{\textbackslash solutiontype} --- Prefix redefinitions for questions and solutions.
    \item \texttt{\textbackslash writtensection} --- Section numbering prefix (e.g. \texttt{5} for 5.1, 5.2).
    \item \texttt{\textbackslash ifcompileanswers} --- Boolean conditional for toggling solution visibility.
\end{itemize}

\section{Code Listings with Minted}
To include syntax-highlighted code in assignments or solutions, use the \texttt{minted} environment (wrap within \texttt{\textbackslash LRE\{\}} when inside a Persian context):

\begin{Verbatim}[frame=lines,framesep=2mm,numbers=left,fontsize=\footnotesize]
import numpy as np

def solve_system(A, b):
    x = np.linalg.solve(A, b)
    residual = np.linalg.norm(np.dot(A, x) - b)
    return x, residual
\end{Verbatim}

\section{Subfiles \& Build Setup}
The template supports modular multi-file assignments with the \texttt{subfiles} package:
\begin{itemize}[leftmargin=1.5em, itemsep=0.2em]
    \item Always set \texttt{\%!TEX root = ../HW.tex} at the top of question files in \texttt{Q-A/}.
    \item In VS Code LaTeX Workshop, keep \texttt{"latex-workshop.latex.rootFile.useSubFile": false} so the master document is always built.
    \item Use \texttt{latexmkrc} for automatic XeLaTeX configuration on Overleaf.
\end{itemize}

\end{document}
